\documentclass[12pt]{article}
\usepackage[utf8]{inputenc}
\usepackage[T1]{fontenc}
\usepackage[margin=1in]{geometry}
\usepackage{setspace}
\usepackage{xcolor}
\usepackage[colorlinks=true, urlcolor=blue!60!black, linkcolor=blue!60!black]{hyperref}
\setlength{\parindent}{0em}
\setlength{\parskip}{0.6em}
\renewcommand{\baselinestretch}{1.0}
\begin{document}
\begin{center}
{\LARGE \textit{Dames \& Moore v. Regan}, 453 U.S. 654 (1981)\par}
\vspace{0.5em}
{\large Case Brief --- Two-Page Edition\par}
\end{center}

\section{Memory Jogger}
To end the Iran Hostage Crisis, the President could nullify attachments on frozen Iranian assets under IEEPA and, through congressional acquiescence, suspend private claims in favor of an international tribunal (\href{https://supreme.justia.com/cases/federal/us/453/654/}{Justia}).

\section{Detailed Case Facts}
After the 1979 embassy seizure, President Carter froze Iranian assets under IEEPA; Treasury licenses let creditors such as Dames \& Moore obtain revocable prejudgment attachments. The 1981 Algiers Accords, which freed the hostages, required the United States to nullify attachments, transfer the assets, and send claims to the Iran--United States Claims Tribunal. Executive orders by Carter and Reagan did so.

\section{Procedural History}
After its attachments were vacated, Dames \& Moore sued Treasury Secretary Regan; the district court dismissed, and the Court granted certiorari before judgment, deciding the case eight days after argument.

\section{Judicial Votes}
Affirmed. Rehnquist wrote for the Court, unanimous on the executive-power judgment; Stevens did not join Part V, and Powell dissented from footnote 6.

\section{Holding}
IEEPA expressly authorized the President to nullify the attachments and transfer the assets, placing him in \textit{Youngstown} Category One. The President could suspend pending claims because claims settlement was necessary to resolve a major dispute over foreign policy, an alternative forum existed, and Congress had long acquiesced. A Tucker Act takings suit remains available.

\section{Analysis of Opinions}
\textbf{Majority:} \textit{Youngstown}'s categories form a spectrum; IEEPA's text covered the attachments, and the ``general tenor'' of IEEPA, the Hostage Act, and the claims-settlement statutes, plus a long practice, supported the suspension. Suspension changed the governing law rather than stripping jurisdiction. \textbf{Stevens:} the takings question was too remote to reach. \textbf{Powell:} attachments are property, and foreign-affairs powers cannot displace the Just Compensation Clause.

\section{Examples: Future Applications}
\textbf{Same-Side:} an executive claims agreement preempts conflicting state law (\textit{Garamendi}). \textbf{Opposite-Side:} a presidential order to state courts enforcing a non-self-executing treaty fails (\textit{Medell\'in}). \textbf{Fence-Sitter:} sweeping IEEPA emergency tariffs, which lack claims settlement's historical gloss.

\section{Critique}
Critics say the Court turned statutory silence into authorization, reading IEEPA---a statute meant to curb emergency powers---as an invitation. Formalists warn that historical gloss creates a one-way ratchet favoring the Executive.

\section{Key Quotations}
\begin{quote}
``Past practice does not, by itself, create power, but `long-continued practice, known to and acquiesced in by Congress, would raise a presumption that the [action] had been [taken] in pursuance of its consent. . . .'\,''
\end{quote}
\noindent --- Justice Rehnquist, for the Court, 453 U.S. at 686.

\end{document}
